\documentclass[12pt]{article}
\usepackage[utf8]{inputenc}
\usepackage[T1]{fontenc}
\usepackage{hyperref}
\title{Sherpa Language Resources for K-12 Curriculum Development}
\author{Compiled by Claude Sonnet 5.5 (claude-sonnet-5-5) for Andr\'as Kornai}
\date{}
\begin{document}
\maketitle

\section{Introduction}
Sherpa (ISO 639-3: xsr; Glottocode sher1255) is a Sino-Tibetan language of Nepal, China and India. This report is based on a single web search. It found reference works and Indian (Sikkim) school materials. It found no Nepali school curriculum for Sherpa, and none of the entries below were opened or checked beyond the search listing.

\section{Primary Educational Resources}
\begin{itemize}
\item CBSE Sherpa secondary syllabus 2024--25. It names the prescribed grammar and writing book \emph{Sherpai Sumtak Tang Tri-Tsom}, which comes from the Sikkim Education Department. This is Indian, not Nepali, school material. \url{https://www.cbseacademic.nic.in/web_material/CurriculumMain25/Sec/Sherpa_Sec_2024-25.pdf}
\item Bharatavani lists Sherpa learning aids, including a grammar and composition book for Classes IX--XII. Publisher and date were not verified. \url{https://bharatavani.in/sherpa/bhashakosha}
\end{itemize}

\section{Linguistic and Reference Materials}
\begin{itemize}
\item \emph{Sherpa-English and English-Sherpa Dictionary with Literary Tibetan and Nepali Equivalents} (Vajra Publications, Kathmandu, 2009), with Tournadre, Lhakpa Norbu Sherpa and others named as authors. \url{https://vedamsbooks.com/books/sherpa_english_and_english_sherpa_dictionary_with_literary_tibetan_and_nepali_equivalents}
\item Barbara Kelly (2004), ``A Grammar and Glossary of the Sherpa Language'', in Genetti (ed.), \emph{Tibeto-Burman Languages of Nepal: Manange and Sherpa}. This is an academic grammar. \url{https://glottolog.org/resource/reference/id/100639}
\item Bharatavani Sherpa dictionary listing, which includes a multilingual dictionary (English--Sherpa--Bhutia--Nepali--Tibetan). \url{https://sec.bharatavani.in/sherpa/bhashakosha}
\end{itemize}

\section{Religious and Community Resources}
The search found nothing specific to Sherpa religious or community organisations.

\section{Digital and Audio Resources}
The search found no Sherpa audio or software resources. The Bharatavani portal is the only digital resource found.

\section{Curriculum Development Context}
A Sherpa-medium curriculum in Nepal would have to be built mostly from scratch. Sikkim's Sherpa syllabus and grammar book and the 2009 dictionary are the closest foundations found. Their suitability for Nepali Sherpa varieties was not assessed.
\end{document}